\section{Related Work}
\label{sec:related}

We review four lines of work related to our problem: UAV relaying and data collection under time constraints, backhaul-aware placement and station selection, the solution methods used for joint trajectory and resource design, and scheduling with release dates and deadlines. \Cref{tab:related} compares the most relevant works along the attributes that matter for deadline-aware alert delivery; its entries are taken from the abstracts of the cited works and from the full text of~\cite{huang2025ddatsap}, and ``Not stated'' marks an attribute that is not mentioned.

% Literature comparison — rows read from data/related_work.csv (review in docs/literature/review.md).
\begin{table}[htbp]
\centering
\footnotesize
\setlength{\tabcolsep}{3pt}
\renewcommand{\arraystretch}{0.92}
\caption{Comparison of related works.}
\label{tab:related}
\newcolumntype{L}[1]{>{\raggedright\arraybackslash}p{#1}}
\begin{tabular}{@{}L{0.06\textwidth}L{0.17\textwidth}L{0.13\textwidth}L{0.14\textwidth}L{0.12\textwidth}L{0.13\textwidth}L{0.15\textwidth}@{}}
\toprule
\textbf{Work} & \textbf{Objective} & \textbf{Station selection, backhaul} & \textbf{Time constraint} & \textbf{Channel, uncertainty} & \textbf{Resources} & \textbf{Method} \\
\midrule
\csvreader[separator=semicolon, late after line=\\, late after last line=\\\midrule]{data/related_work.csv}{1=\colkey, 4=\colobjective, 6=\colrelay, 7=\coldeadline, 9=\colchannel, 10=\colresource, 11=\colmethod}{\expandafter\cite\expandafter{\colkey} & \colobjective & \colrelay & \coldeadline & \colchannel & \colresource & \colmethod}
This work & Timely-delivery ratio & Yes, ground relays with finite backhaul & Yes, per-alert deadline & PrLoS, design on sampled realizations & Weighted bandwidth & Evaluated block search, schedule repair \\
\bottomrule
\end{tabular}
\end{table}


Huang et al.~\cite{huang2025ddatsap} use a UAV as a buffered decode-and-forward relay for multiple user pairs under the PrLoS channel and maximize the minimum average rate through BCD combined with SCA; our work inherits their channel model and three-block structure. The works closest to timely alert delivery maximize the number of devices served. Tran et al.~\cite{tran2022relay} require the data of each device to reach a ground gateway in time via the UAV's memory and optimize bandwidth, power, and trajectory through binary-variable relaxation and inner approximation. Samir et al.~\cite{samir2020time} assign a hard deadline to each device, obtain the global optimum with branch, reduce and bound for small instances, and resort to SCA for large ones, while Du et al.~\cite{du2023timeconstrained} count the devices served within a finite collection time. Other works place time in the constraints or the objective in different ways: a common collection deadline with an MILP model and a heuristic~\cite{albusalih2018deadline}, an age-of-information threshold together with the selection of the ground station that receives the data~\cite{liu2022aoi,liu2024interference}, bounded queueing delay in an online Lyapunov-based design~\cite{tang2022delay}, delay as the objective of UAV relaying or of priority-based collection~\cite{he2022delay,cao2023priority,fadlullah2016dynamic}, and message time-to-live in a delay-tolerant network with reinforcement-learning UAVs~\cite{wang2026juror}.

When a UAV acts as a base station or relay connecting users to the core network, the backhaul capacity limits the number of users that can be served. Kalantari et al.~\cite{kalantari2017backhaul} place flying base stations in 3D according to the backhaul type; Selim and Kamal~\cite{selim2018postdisaster} employ tethered drones as post-disaster backhaul; Yu et al.~\cite{yu2023backhaul} jointly optimize bandwidth and placement under a free-space-optics backhaul capacity constraint; and Queiros et al.~\cite{queiros2024predictive} allocate bandwidth and position second-tier flying relays given the known trajectories of the first tier. None of these works imposes per-request time constraints.

BCD combined with SCA is the prevailing approach to joint trajectory and resource design~\cite{huang2025ddatsap,wu2018multi,liu2024interference}; Wu et al.~\cite{wu2018multi} prove that it converges to at least a locally optimal solution of the max--min throughput problem. Works that count served requests must handle binary variables, either by relaxation~\cite{tran2022relay} or by branch-and-bound~\cite{samir2020time}. Park and Lee~\cite{park2026expected} show that approximating the expected rate by the rate at the average channel overestimates the true rate because of Jensen's inequality, which motivates our choice to score designs on sampled realizations. Large neighborhood search with destroy and repair operators~\cite{ropke2006alns} is a precedent for our repair mechanism, and routing over time-varying graphs with known dynamics~\cite{jain2004dtn} is a precedent for our time-respecting connectivity measure.

On a single link, selecting the set of alerts delivered on time resembles single-machine scheduling to minimize the number of late jobs. The non-preemptive version $1\,|\,r_j\,|\,\sum U_j$ is strongly NP-hard~\cite{lenstra2020elements}, whereas the preemptive version $1\,|\,r_j,\mathrm{pmtn}\,|\,\sum U_j$ is solvable in polynomial time by dynamic programming~\cite{lawler1990preemptive}. Our model allows data to be split across slots, which is closer to the preemptive case; as we show, its hardness stems instead from the requirement that each alert follow exactly one path. Our upper bound relies on the concavity of the rate in the bandwidth and on its envelope of tangents, a standard tool of convex optimization~\cite{boyd2004convex}.

As \cref{tab:related} shows, no existing work combines per-request deadlines with the selection of ground relays over a backhaul of finite capacity: works with deadlines assume a single destination that accepts all data, and works with backhaul constraints have no per-request deadlines. Moreover, the methods that count served requests design on a deterministic channel, whereas the PrLoS channel makes the delivered count a random quantity. This paper addresses both gaps. As the strongest available baseline, we adopt the half-duplex algorithm of Tran et al.~\cite{tran2022relay}: its objective is the closest to the number of timely alerts, it designs offline for known demand and jointly optimizes trajectory and radio resources, and it is the only one that models data passing through the UAV buffer to a ground gateway, which corresponds to our paths on which the UAV drops data into the ground network.
